\chapter{Game behavior}
\label{g:behavior}

\section{Core loop}
Il loop (riassunto in Fig.~\ref{fig:coreloop}) procede cos\`i:
\begin{enumerate}
\item \textbf{Studio}: pagina di dettaglio del topic (testo breve + link
opzionali). Nessun timer.
\item \textbf{Quiz}: 5 domande a risposta multipla, soglia di superamento
\textbf{80\,\% (4/5)}. I tentativi sono illimitati ma penalizzati (nessuna
XP in caso di sconfitta/ripetizione fallita, perdita di vite nel boss).
\item \textbf{Ricompensa}: a quiz superato, scelta \textbf{1-di-3} carte
(per tipo di ricompensa). Una sola ricompensa per topic; inventario visibile.
\item \textbf{Avanzamento}: lo sblocco del nodo successivo lungo la roadmap;
a fine capitolo, \textbf{boss fight}.
\end{enumerate}

\begin{figure}[ht]
\centering
\resizebox{\linewidth}{!}{%
\begin{tikzpicture}[
  node distance=8mm,
  every node/.style={font=\scriptsize},
  cstep/.style={draw, rounded corners=3pt, fill=blue!8, minimum width=2.4cm, minimum height=9mm, align=center},
  >=Stealth
]
\node[cstep] (s) {Studio\\topic};
\node[cstep, right=of s] (q) {Quiz 5Q\\soglia 80\%};
\node[cstep, right=of q] (r) {Ricompensa\\1-di-3};
\node[cstep, right=of r] (a) {Avanzamento\\sblocco nodo};
\node[cstep, right=of a] (b) {Boss fight\\10 HP vs 3 HP};
\draw[->] (s) -- (q); \draw[->] (q) -- (r); \draw[->] (r) -- (a); \draw[->] (a) -- (b);
\draw[->, dashed] (b.south) -- ++(0,-10mm) -| (s.south);
\node[anchor=north, font=\scriptsize\itshape] at ([yshift=-11mm]q.south) {vittoria boss: capitolo successivo, si ritorna a Studio};
\end{tikzpicture}}
\caption{Core loop: Studio $\rightarrow$ Quiz $\rightarrow$ Ricompensa $\rightarrow$ Avanzamento $\rightarrow$ Boss fight, con ritorno al capitolo successivo.}
\label{fig:coreloop}
\end{figure}

\section{Boss fight}
Combattimento a turni con numeri reali (stato vero implementato):
\begin{itemize}
\item Boss: \textbf{10 HP}; giocatore: \textbf{3 HP}; \textbf{vite totali 3}.
\item \textbf{Energia 3 per turno} (1 per carta giocata).
\item A fine turno quiz obbligatorio: risposta giusta = danno al boss
(quiz-danno), risposta errata = danno al giocatore.
\item Vittoria: claim delle ricompense del nodo boss (semplificazione
dichiarata: tutte quelle del nodo, non 1 sola), sblocco del capitolo
successivo, badge del boss, \textbf{+100 XP} (come l'Hub). Sconfitta:
retry a HP pieni, senza XP.
\end{itemize}
I tre boss (con lore di una riga dalla campagna Web):
\begin{itemize}
\item \textbf{Man-in-the-Middle} (cap.~1): ``vive nei Wi-Fi aperti e legge
ci\`o che non \`e cifrato'';
\item \textbf{The Amnesiac} (cap.~2): ``ha cancellato la memoria del web: ti
costringe a ricordare tutto da solo, tra bigliettini, sessioni e copie
stantie'';
\item \textbf{Spaghetti Colossus} (cap.~3): ``un ammasso di pagine copiate e
mai versionate: ogni modifica spezza qualcos'altro''.
\end{itemize}
I quiz del boss sono adattivi: pescati dai 3 topic del capitolo.

\section{Economie numeriche locali}
\begin{itemize}
\item \textbf{XP/livelli}: soglie \textbf{L1 0 / L2 100 / L3 500}, identiche
a quelle dell'Hub cos\`i che locale e remoto coincidano. Quiz superato:
\textbf{+100 XP}; vittoria boss: \textbf{+100 XP} (stessi dell'Hub).
\item \textbf{Streak}: conteggiata ai fini del ripasso
(\path{failCount}, soglia di ripasso); il bonus XP \`e definito ma
\textbf{non attivo} e non va considerato nella valutazione.
\item \textbf{Vite}: 3. \textbf{Save per utente}: SharedPreferences
(\texttt{Save v1}: nodi completati, deck, XP, livello, boss sbloccati,
reward riscattate, badge; \texttt{Continue} ripristina,
\texttt{New Run} chiede conferma e azzera solo quell'utente).
Auth locale con hash salato e migrazione dalle password storiche.
\item \textbf{Animazioni}: sobrie (transizioni essenziali, niente effetti
extra).
\end{itemize}
